% ============================================================================
%  Chapitre 2, partie IV : Mécanique analytique
%  (manuscrit P2, p. 11 à 24, et P3, p. 1 à 2)
% ============================================================================

% ----------------------------------------------------------------------------
\section{Mécanique analytique}
% ----------------------------------------------------------------------------

% ----------------------------------------------------------------------------
\subsection{Forces généralisées et travail virtuel}

\subsubsection{Déplacements virtuels et travail virtuel}

\paragraph{Variation de la position d'un point}

Soit un système décrit localement par les coordonnées
$q = (q_1,\dots,q_n)$. La position d'un point matériel $M_a$ s'écrit
$\vect{r}_a = \vect{r}_a(q,t)$. Pour une variation de configuration à
temps fixé, $\delta t = 0$, sa variation de position est
\[
    \delta\vect{r}_a = \sum_{i=1}^{n}\pd{\vect{r}_a}{q_i}\,\delta q_i
\]
En notant $\mat{J}_a$ la jacobienne de $\vect{r}_a$ par rapport à $q$,
\[
    \mat{J}_a(q,t) = \begin{bmatrix}
        \pd{\vect{r}_a}{q_1} & \cdots & \pd{\vect{r}_a}{q_n}
    \end{bmatrix},
    \qquad
    \delta\vect{r}_a = \mat{J}_a(q,t)\,\delta q
\]

\paragraph{Travail virtuel}

\begin{definition}[Travail virtuel d'une force]
    Le \emph{travail virtuel} d'une force $\vect{F}_a$ appliquée au
    point $M_a$ est
    \[
        \delta W_a = \vect{F}_a\cdot\delta\vect{r}_a
    \]
    Pour $N$ forces, $\delta W = \sum_{a=1}^{N}\vect{F}_a\cdot
    \delta\vect{r}_a$.
\end{definition}

\subsubsection{Forces généralisées}

En remplaçant $\delta\vect{r}_a$ et en échangeant les sommes,
\[
    \delta W
    = \sum_{a=1}^{N}\vect{F}_a\cdot
    \left(\sum_{i=1}^{n}\pd{\vect{r}_a}{q_i}\,\delta q_i\right)
    = \sum_{i=1}^{n}\left(\sum_{a=1}^{N}\vect{F}_a\cdot
    \pd{\vect{r}_a}{q_i}\right)\delta q_i
\]

\begin{definition}[Force généralisée]
    La \emph{force généralisée} associée à la coordonnée $q_i$ est le
    coefficient $Q_i$ de $\delta q_i$ dans le travail virtuel :
    \[
        Q_i = \sum_{a=1}^{N}\vect{F}_a\cdot\pd{\vect{r}_a}{q_i},
        \qquad
        \vect{Q} = \sum_{a=1}^{N}\mat{J}_a^T\vect{F}_a
    \]
    Alors $\delta W = \sum_{i=1}^{n}Q_i\,\delta q_i = \vect{Q}^T\delta q$.
\end{definition}

\subsubsection{Forces dérivant d'un potentiel}

\begin{definition}[Forces dérivant d'un potentiel]
    Des forces \emph{dérivent d'une énergie potentielle} $V(q,t)$
    lorsque leur travail virtuel vérifie, à temps fixé,
    \[
        \delta W_{\mathrm{pot}} = -\,\delta V
    \]
\end{definition}

Comme $\delta V = \sum_i\pd{V}{q_i}\,\delta q_i$, on identifie
\[
    Q_i^{\mathrm{pot}} = -\pd{V}{q_i}
\]
Les autres forces appliquées sont dites \emph{non conservatives}, de
forces généralisées $\vect{Q}^{\mathrm{nc}}$.

\subsubsection{Puissance en coordonnées généralisées}

Si $\vect{r}_a = \vect{r}_a(q)$ ne dépend pas du temps,
$\dot{\vect{r}}_a = \mat{J}_a\dot q$ et la puissance totale est
\[
    \mathcal{P} = \sum_{a=1}^{N}\vect{F}_a\cdot\dot{\vect{r}}_a
    = \underbrace{\sum_{a=1}^{N}\vect{F}_a^T\mat{J}_a}_{\vect{Q}^T}\,\dot q
    = \vect{Q}^T\dot q
\]
Si $\vect{r}_a = \vect{r}_a(q,t)$, alors $\dot{\vect{r}}_a =
\mat{J}_a\dot q + \pd{\vect{r}_a}{t}$ et
\[
    \mathcal{P} = \vect{Q}^T\dot q
    + \sum_{a=1}^{N}\vect{F}_a^T\pd{\vect{r}_a}{t}
\]

% ----------------------------------------------------------------------------
\subsection{Action et Principe de Hamilton}

On considère un système à contraintes holonomes idéales, décrit par
les coordonnées généralisées $q$, dont les forces appliquées dérivent
d'une énergie potentielle $V(q,t)$.

\subsubsection{Le lagrangien}

\begin{definition}[Lagrangien du système]
    Le \emph{lagrangien} du système est
    \[
        L(q,\dot q,t) = T(q,\dot q,t) - V(q,t)
    \]
    où $T$ est l'énergie cinétique et $V$ l'énergie potentielle.
\end{definition}

\subsubsection{L'action}

\begin{definition}[Action]
    L'\emph{action} associée à une trajectoire $q : [t_0,t_1]\to Q$
    est
    \[
        S[q] = \int_{t_0}^{t_1} L\bigl(q(t),\dot q(t),t\bigr)\,\dd t
    \]
\end{definition}

L'action est une fonctionnelle : elle associe un nombre à une
trajectoire entière. Son unité est le $\mathrm{J\cdot s}$ (joule-seconde). Pour comparer une trajectoire à ses voisines, on introduit
\[
    q_\varepsilon(t) = q(t) + \varepsilon\,\eta(t),
    \qquad \eta(t_0) = \eta(t_1) = 0
\]
On fixe les extrêmités donc les perturbations y sont nulles.

\subsubsection{Principe de moindre action}

L'idée de départ est formulée par Maupertuis (Mémoires de l'Académie de Sciences de paris du 15 avril 1744) :
\begin{quote}
    \itshape « Lorsqu'il arrive quelque changement dans la Nature, la
    quantité d'action nécessaire pour ce changement est la plus petite
    qu'il soit possible. »
\end{quote}
Traduisons alors cette idée. Supposons que la trajectoire physique $q^*$
minimise localement l'action $S$. Pour toute perturbation $\eta$
s'annulant aux bords, on doit avoir
\[
    S[q^* + \varepsilon\eta] \geq S[q^*]
\]
pour tout $\varepsilon$ suffisamment petit. La fonction réelle
$a(\varepsilon) = S[q^* + \varepsilon\eta]$ possède alors un minimum
en $\varepsilon = 0$, donc $a'(0) = 0$. Une trajectoire qui minimise
l'action satisfait ainsi
\[
    \delta S[q^*;\eta] = 0
    \qquad \forall\,\eta,\ \eta(t_0) = \eta(t_1) = 0
\]

\subsubsection{Principe de Hamilton}
\begin{theoreme}[Principe de Hamilton]
    La trajectoire physique rend l'action stationnaire parmi les
    trajectoires voisines ayant les mêmes extrémités :
    \[
        \delta S[q;\eta] := \frac{\dd}{\dd\varepsilon}
        S[q + \varepsilon\eta]\Big|_{\varepsilon=0} = 0
    \]
    pour toute variation $\eta$ suffisamment régulière s'annulant aux
    extrémités.
\end{theoreme}

\begin{note}
    Si on déforme légèrement la trajectoire physique en gardant les
    extrémités fixées, l'action ne change pas au premier ordre. La
    trajectoire n'est pas forcément un minimum : ce n'est vrai que sur
    des intervalles de temps assez courts, d'où le terme
    \emph{stationnaire}.
\end{note}

% ----------------------------------------------------------------------------
\subsection{Équations d'Euler-Lagrange}

\subsubsection{Calcul de la première variation}

On effectue ici les mêmes calculs qu'au chapitre 1.

\begin{proposition}[Première variation de l'action]
    Si $L$ est de classe $C^2$, la trajectoire $q$ de classe $C^2$ et
    $\eta$ de classe $C^1$ avec $\eta(t_0) = \eta(t_1) = 0$, alors
    \[
        \delta S[q;\eta] = \int_{t_0}^{t_1}\sum_{i=1}^{n}
        \left(\pd{L}{q_i} - \ddt\left(\pd{L}{\dot q_i}\right)\right)
        \eta_i\,\dd t
    \]
\end{proposition}

\begin{demonstration}
    On a $q_\varepsilon(t) = q(t) + \varepsilon\eta(t)$ avec
    $\eta(t_0) = \eta(t_1) = 0$, donc $\dot q_\varepsilon(t) =
    \dot q(t) + \varepsilon\dot\eta(t)$. La fonctionnelle associée à
    $q_\varepsilon$ est
    \[
        S[q_\varepsilon] = \int_{t_0}^{t_1}
        L(q+\varepsilon\eta,\dot q+\varepsilon\dot\eta,t)\,\dd t
    \]
    On calcule alors la première variation en dérivant sous
    l'intégrale :
    \[
        \delta S[q;\eta]
        = \frac{\dd}{\dd\varepsilon}\left(\int_{t_0}^{t_1}
        L(q+\varepsilon\eta,\dot q+\varepsilon\dot\eta,t)\,\dd t\right)
        \bigg|_{\varepsilon=0}
        = \int_{t_0}^{t_1}\left(\frac{\dd}{\dd\varepsilon}
        L(q+\varepsilon\eta,\dot q+\varepsilon\dot\eta,t)
        \Big|_{\varepsilon=0}\right)\dd t
    \]
    On considère l'intégrande. Par la règle de chaîne, \textbf{en postulant
    $q$ et $\dot q$ indépendants},
    \begin{align*}
        \frac{\dd}{\dd\varepsilon}
        L(q+\varepsilon\eta,\dot q+\varepsilon\dot\eta,t)
        \Big|_{\varepsilon=0}
        &= \left(\pd{L}{(q+\varepsilon\eta)}\,
        \pd{(q+\varepsilon\eta)}{\varepsilon}
        + \pd{L}{(\dot q+\varepsilon\dot\eta)}\,
        \pd{(\dot q+\varepsilon\dot\eta)}{\varepsilon}\right)
        \bigg|_{\varepsilon=0}\\
        &= \pd{L}{q}\,\eta + \pd{L}{\dot q}\,\dot\eta
    \end{align*}
    Donc, en écrivant les composantes,
    \[
        \delta S[q;\eta] = \int_{t_0}^{t_1}\sum_{i=1}^{n}
        \left(\pd{L}{q_i}\,\eta_i + \pd{L}{\dot q_i}\,\dot\eta_i\right)
        \dd t
    \]
    On considère ensuite $I := \int_{t_0}^{t_1}\pd{L}{\dot q_i}\,
    \dot\eta_i\,\dd t$ et on intègre par parties avec
    \[
        u = \pd{L}{\dot q_i} \longrightarrow
        \dd u = \ddt\left(\pd{L}{\dot q_i}\right)\dd t,
        \qquad
        \dd v = \dot\eta_i\,\dd t \longrightarrow v = \eta_i
    \]
    Comme $\eta_i(t_0) = \eta_i(t_1) = 0$, le terme de bord est nul :
    \[
        I = \underbrace{\left[\pd{L}{\dot q_i}\,\eta_i
        \right]_{t_0}^{t_1}}_{=\,0}
        - \int_{t_0}^{t_1}\ddt\left(\pd{L}{\dot q_i}\right)\eta_i\,
        \dd t
        = -\int_{t_0}^{t_1}\ddt\left(\pd{L}{\dot q_i}\right)\eta_i\,
        \dd t
    \]
    En le réinjectant dans $\delta S$, on obtient le résultat.
\end{demonstration}

\subsubsection{Équations d'Euler-Lagrange}

Si $\delta S = 0$ pour tout $\eta$, le lemme de Du Bois-Reymond,
appliqué coordonnée par coordonnée avec $\eta\in C^1$ nulle aux bords,
donne
\[
    \pd{L}{q_i} - \ddt\left(\pd{L}{\dot q_i}\right) = 0
\]

\begin{theoreme}[Équations d'Euler-Lagrange]
    Sous les hypothèses de régularité et les conditions aux limites
    précédentes, $q$ rend l'action stationnaire si et seulement si $q$
    satisfait les \emph{équations d'Euler-Lagrange} :
    \[
        \pd{L}{q_i} - \ddt\left(\pd{L}{\dot q_i}\right) = 0,
        \qquad i = 1,\dots,n
    \]
\end{theoreme}

% ----------------------------------------------------------------------------
\subsection{Principe de d'Alembert et dynamique non holonome}

On introduit ce principe pour décrire les mouvements sous
contraintes.

\subsubsection{Principe de d'Alembert}

On considère des points matériels de masses constantes $m_a$ et de
positions $\vect{r}_a$. On note $\vect{F}_a$ les forces appliquées et
$\vect{R}_a$ les réactions de liaison. La deuxième loi de Newton donne
\[
    m_a\ddot{\vect{r}}_a = \vect{F}_a + \vect{R}_a
\]
Pour un déplacement virtuel admissible $\delta\vect{r}_a$, on
multiplie scalairement puis on somme :
\[
    \sum_a(\vect{F}_a - m_a\ddot{\vect{r}}_a)\cdot\delta\vect{r}_a
    + \sum_a\vect{R}_a\cdot\delta\vect{r}_a = 0
\]

\begin{hypotheses}[Liaisons idéales]
    Les réactions ont un travail virtuel total nul pour tous les
    déplacements virtuels admissibles :
    \[
        \sum_a\vect{R}_a\cdot\delta\vect{r}_a = 0
    \]
\end{hypotheses}

\begin{theoreme}[Principe de d'Alembert]
    Pour des liaisons idéales, pour tout déplacement virtuel
    admissible,
    \[
        \sum_a(\vect{F}_a - m_a\ddot{\vect{r}}_a)\cdot
        \delta\vect{r}_a = 0
    \]
\end{theoreme}

\subsubsection{Contraintes cinématiques et déplacements virtuels admissibles}

Les contraintes utilisées sont $\mat{A}(q)\,\dot q = 0$, avec
$\mat{A}(q)\in\mathcal{M}_{k,n}(\R)$ de rang $k$. Dans le , \textbf{les déplacements virtuels} satisfont
\[
    \underbrace{\mat{A}(q)\,\delta q = 0}
    _{\text{restriction des déplacements virtuels}}
    \qquad\text{et non}\qquad
    \underbrace{\mat{A}(q)\,\dot q = 0}
    _{\text{contrainte sur le mouvement}}
\]

\subsubsection{De d'Alembert à la formulation lagrangienne}

On considère une variation $\eta(t) = \delta q(t)$ admissible à chaque
instant et nulle aux extrémités :
\[
    \mat{A}\bigl(q(t)\bigr)\,\eta(t) = 0,
    \qquad \eta(t_0) = \eta(t_1) = 0
\]
L'énergie cinétique du système est $T = \frac12\sum_a m_a
\|\dot{\vect{r}}_a\|^2$, donc à temps fixé
\[
    \delta T = \sum_a m_a\dot{\vect{r}}_a\cdot\delta\dot{\vect{r}}_a
    = \sum_a m_a\dot{\vect{r}}_a\cdot\ddt\delta\vect{r}_a
\]
En intégrant par parties, le terme de bord est nul car les variations
s'annulent aux extrémités :
\[
    \int_{t_0}^{t_1}\sum_a m_a\ddot{\vect{r}}_a\cdot\delta\vect{r}_a
    \,\dd t
    = \Bigl[\sum_a m_a\dot{\vect{r}}_a\cdot\delta\vect{r}_a
    \Bigr]_{t_0}^{t_1}
    - \int_{t_0}^{t_1}\delta T\,\dd t
    = -\,\delta\int_{t_0}^{t_1}T\,\dd t
\]
Par ailleurs, en séparant les forces dérivant d'un potentiel des
forces non conservatives,
\[
    \sum_a\vect{F}_a\cdot\delta\vect{r}_a
    = -\,\delta V + (\vect{Q}^{\mathrm{nc}})^T\eta
\]
Le principe de d'Alembert, intégré entre $t_0$ et $t_1$, devient
\begin{align*}
    0 &= \sum_a (\vect{F}_a - m_a \ddot{\vect{r}}_a) \cdot \delta \vect{r}_a \\
    \leftrightarrow 0 &= \int_{t_0}^{t_1}\Bigl(-\,\delta V
    + (\vect{Q}^{\mathrm{nc}})^T\eta\Bigr)\dd t
    - \int_{t_0}^{t_1}\sum_a m_a\ddot{\vect{r}}_a\cdot
    \delta\vect{r}_a\,\dd t\\
    &= \int_{t_0}^{t_1}\delta(T - V)\,\dd t
    + \int_{t_0}^{t_1}(\vect{Q}^{\mathrm{nc}})^T\eta\,\dd t
    = \delta\int_{t_0}^{t_1}L\,\dd t
    + \int_{t_0}^{t_1}(\vect{Q}^{\mathrm{nc}})^T\eta\,\dd t
\end{align*}

\begin{theoreme}[Principe de Lagrange-d'Alembert]
    Le mouvement vérifie
    \[
        \delta S[q;\eta]
        + \int_{t_0}^{t_1}(\vect{Q}^{\mathrm{nc}})^T\eta\,\dd t = 0
    \]
    pour toutes les variations telles que $\mat{A}(q)\,\eta = 0$ et
    $\eta(t_0) = \eta(t_1) = 0$, et il satisfait
    $\mat{A}(q)\,\dot q = 0$.
\end{theoreme}

\subsubsection{Équations d'Euler-Lagrange-d'Alembert}

On introduit
\[
    \mathcal{L}(L) = \ddt\pd{L}{\dot q} - \pd{L}{q},
    \qquad\text{de sorte que}\qquad
    \delta S[q;\eta] = -\int_{t_0}^{t_1}\mathcal{L}(L)^T\eta\,\dd t
\]
Le principe de Lagrange-d'Alembert devient
\[
    \int_{t_0}^{t_1}\bigl(\mathcal{L}(L) - \vect{Q}^{\mathrm{nc}}\bigr)^T
    \eta\,\dd t = 0
\]
Le lemme de Du Bois-Reymond, avec les mêmes hypothèses sur $\eta$,
donne
\[
    \bigl(\mathcal{L}(L) - \vect{Q}^{\mathrm{nc}}\bigr)^T\delta q = 0
    \qquad \forall\,\delta q\in\Ker\mat{A}(q)
\]
car $\mat{A}(q)\,\eta = 0 \leftrightarrow \eta\in\Ker\mat{A}(q)$. On
sait seulement que $\mathcal{L}(L) - \vect{Q}^{\mathrm{nc}}$, la réaction
généralisée nécessaire pour respecter les contraintes, est orthogonale
à $\Ker\mat{A}(q)$. Par l'orthogonalité noyau-image,
\[
    \mathcal{L}(L) - \vect{Q}^{\mathrm{nc}}\in(\Ker\mat{A})^\perp
    = \Ima\mat{A}^T
\]
Il existe donc $\vect{\lambda}\in\R^k$ tel que
$\mathcal{L}(L) - \vect{Q}^{\mathrm{nc}} = \mat{A}(q)^T\vect{\lambda}$.

\begin{theoreme}[Équations d'Euler-Lagrange-d'Alembert]
    Le mouvement est solution de
    \[
        \left\{
        \begin{aligned}
            &\ddt\pd{L}{\dot q} - \pd{L}{q}
            = \vect{Q}^{\mathrm{nc}} + \mat{A}(q)^T\vect{\lambda},
            \qquad \vect{\lambda}\in\R^k\\
            &\mat{A}(q)\,\dot q = 0
        \end{aligned}
        \right.
    \]
    où $k$ est le nombre de contraintes et $\mat{A}(q)^T\vect{\lambda}$ les
    réactions généralisées.
\end{theoreme}

\subsubsection{Exemple : roue sur ressort et force constante}
On va faire ici une analyse complète sur l'application de l'équation Euler-Lagrange-d'Alembert sur un modèle proche de ce que l'on va travailler en terme de contraintes.

Une roue (disque) de masse $m$, de rayon $r$ et de moment d'inertie
$J$ autour de son axe roule sans glisser sur un sol horizontal. Son
centre est relié à un mur par un ressort de raideur $k$, et on lui
applique une force horizontale constante $\vect{F} = (F,0)$. On prend
$q = (x,\varphi)$, où $x$ est la position du centre mesurée depuis la
longueur à vide du ressort et $\varphi$ l'angle de rotation de la
roue (\cref{fig-roue-ressort}).

\begin{figure}[H]
    \centering
    % ============================================================================
%  Figure : roue roulant sans glisser, reliée à un mur par un ressort,
%  soumise à une force horizontale constante F
% ============================================================================
\begin{tikzpicture}[>=Stealth, line cap=round]
    % paramètres
    \def\R{1}       % rayon de la roue
    \def\xc{5.2}    % abscisse du centre

    % sol et mur hachurés
    \fill[pattern=north east lines] (-0.3,-0.25) rectangle (8.3,0);
    \fill[pattern=north east lines] (-0.3,0) rectangle (0,2.6);
    \draw[thick] (8.3,0) -- (0,0) -- (0,2.6);

    % ressort du mur au centre
    \draw[thick] (0,\R) -- (0.5,\R);
    \draw[thick, decorate, decoration={coil, aspect=0.5,
        segment length=6pt, amplitude=6pt}]
        (0.5,\R) -- (\xc-1.4,\R);
    \node[above] at (2.2,\R+0.25) {$k$};

    % roue
    \draw[thick, fill=accent!8] (\xc,\R) circle (\R);
    \draw[thick] (\xc-1.4,\R) -- (\xc,\R);   % attache du ressort en C
    \fill (\xc,\R) circle (1.5pt);
    \node[below left] at (\xc,\R) {$C$};
    \draw (\xc,\R) -- ++(60:\R);   % rayon repère de l'angle

    % angle phi
    \draw[->, accent] (\xc,\R) ++(120:0.55) arc (120:30:0.55);
    \node[accent] at ($(\xc,\R)+(80:0.8)$) {$\varphi$};

    % force F au centre
    \draw[->, very thick, couldef] (\xc,\R) -- ++(1.6,0)
        node[right] {$\vect{F}$};

    % point de contact et frottement
    \fill (\xc,0) circle (1.5pt);
    \node[below=2pt] at (\xc,-0.25) {$I$};
    \draw[->, very thick, coulthm] (\xc,0.06) -- ++(-1.1,0)
        node[above] {$\lambda$};

    % rayon r
    \draw[<->] (\xc,0.05) -- (\xc,\R-0.05)
        node[midway, right] {$r$};

    % repérage de x depuis la position à vide du ressort
    \def\xo{3.4}
    \draw[densely dashed, gray] (\xo,-0.25) -- (\xo,2.3);
    \node[gray, above] at (\xo,2.3) {$x = 0$};
    \draw[|->] (\xo,-0.95) -- (\xc,-0.95) node[midway, below] {$x$};
\end{tikzpicture}

    \caption{Roue roulant sans glisser, reliée au mur par un ressort
    et poussée par une force constante $\vect{F}$ ; $\lambda_x$ est le
    frottement au point de contact $I$.}
    \label{fig-roue-ressort}
\end{figure}

\paragraph{Lagrangien, contrainte et force non conservative}

\[
    L(q,\dot q) = \frac12\,m\,\dot x^2 + \frac12\,J\,\dot\varphi^2
    - \frac12\,k\,x^2
\]
On écrit maintenant la contrainte de roulement sans glissement. On
munit $\rep{s}$ de la base $(\vect{e}_x,\vect{e}_y,\vect{e}_z)$, avec
$\vect{e}_x$ horizontal vers la droite, $\vect{e}_y$ vertical vers le
haut et $\vect{e}_z$ sortant de la figure, et d'origine $O$ au sol
sous la position $x = 0$. Le centre reste à la hauteur $r$, d'où $\overrightarrow{OC} = x\,\vect{e}_x + r\,\vect{e}_y$
et
\[
    \vect{v}_C = \dot x\,\vect{e}_x
\]
Le disque tourne autour de l'axe $\vect{e}_z$, qui est fixe à la fois dans
$\rep{s}$ et dans $\rep{b}$. Comme $\mat{R}$ est une rotation d'axe
$\vect{e}_z$, $\mat{R}\,\vect{e}_z = \vect{e}_z$, et avec $+ \varphi$ dans le sens horaire,
\[
    \vect{\omega} = -\,\dot\varphi\,\vect{e}_z,
    \qquad
    \vect{\Omega} = \mat{R}\,\vect{\omega} = -\,\dot\varphi\,\vect{e}_z
\]
Le point de contact $I$ vérifie $\overrightarrow{CI} = -\,r\,\vect{e}_y$.
Le champ des vitesses du solide donne
\[
    \vect{v}_I = \vect{v}_C + \vect{\Omega}\wedge\overrightarrow{CI}
\]
avec
\[
    \vect{\Omega}\wedge\overrightarrow{CI}
    = (-\,\dot\varphi\,\vect{e}_z)\wedge(-\,r\,\vect{e}_y)
    = r\,\dot\varphi\,(\vect{e}_z\wedge\vect{e}_y)
    = -\,r\,\dot\varphi\,\vect{e}_x
\]
d'où
\[
    \vect{v}_I = (\dot x - r\,\dot\varphi)\,\vect{e}_x
\]
En composantes dans $\rep{s}$, la vitesse du point de contact est
linéaire en $\dot q$ : on l'écrit avec la \emph{matrice de contrainte
locale au sol} $\mat{A}_I(q)\in\mathcal{M}_{3,2}(\R)$,
\[
    \vect{v}_I
    = \begin{pmatrix} \dot x - r\,\dot\varphi\\ 0\\ 0 \end{pmatrix}
    = \underbrace{\begin{pmatrix}
        1 & -r\\ 0 & 0\\ 0 & 0
    \end{pmatrix}}_{\mat{A}_I}
    \begin{pmatrix} \dot x\\ \dot\varphi \end{pmatrix}
\]
Le roulement sans glissement impose $\vect{v}_I = \vect{0}$, soit
$\mat{A}_I\,\dot q = 0$. Les deux dernières lignes de $\mat{A}_I$ sont
nulles (le centre reste à hauteur $r$ et le mouvement est plan) :
$\rang\mat{A}_I = 1$ et il ne reste que $k = 1$ contrainte,
\[
    \dot x - r\,\dot\varphi = 0
    \longleftrightarrow
    \mat{A}(q)\,\dot q = 0,
    \qquad
    \mat{A} = \begin{pmatrix} 1 & -r \end{pmatrix}
\]

\begin{remarque}[Nombre de coordonnées généralisées]
    Le théorème du rang justifie le choix $q = (x,\varphi)$. On part
    d'un solide plan libre, décrit par $(x_C,y_C,\varphi)$, soit
    $3$ degrés de liberté, et on examine la contrainte de non-glissement :
    \begin{itemize}
        \item \emph{Contact avec le sol.} Le champ des vitesses donne
            \[
                \vect{v}_I
                = \begin{pmatrix}
                    \dot x_C - r\,\dot\varphi\\ \dot y_C
                \end{pmatrix}
                = \begin{pmatrix}
                    1 & 0 & -r\\ 0 & 1 & 0
                \end{pmatrix}
                \begin{pmatrix}
                    \dot x_C\\ \dot y_C\\ \dot\varphi
                \end{pmatrix}
            \]
            matrice de rang $2$. La ligne $\dot y_C = 0$ peut s'intégrer en
            $y_C = r$ : c'est une contrainte holonome de rang $1$, qui
            retire une coordonnée par définition, d'où $3 - 1 = 2$ coordonnées
            généralisées, $(x,\varphi)$. La ligne restante est la
            contrainte de roulement $\mat{A}\,\dot q = 0$, de rang $1$
            : les vitesses admissibles forment un espace de dimension
            $2 - 1 = 1$.
    \end{itemize}
\end{remarque}

La force $\vect{F}$ est appliquée au centre, donc
\[
    \vect{Q}^{\mathrm{nc}} = \begin{pmatrix} F\\ 0 \end{pmatrix}
\]

\paragraph{Réactions généralisées}

On part de la contrainte locale $\mat{A}_I\,\dot q = 0$ : la liaison
étant idéale, ses réactions généralisées sont de la forme
$\mat{A}_I^T\vect{\lambda}$, avec un multiplicateur par ligne,
$\vect{\lambda} = (\lambda_x,\lambda_y,\lambda_z)\in\R^3$ :
\[
    \vect{Q}^r = \mat{A}_I^T\vect{\lambda}
    = \begin{pmatrix}
        1 & 0 & 0\\ -r & 0 & 0
    \end{pmatrix}
    \begin{pmatrix} \lambda_x\\ \lambda_y\\ \lambda_z \end{pmatrix}
    = \begin{pmatrix} \lambda_x\\ -\,r\,\lambda_x \end{pmatrix}
\]
Ces multiplicateurs ont un sens physique. Comme
$\vect{v}_I = \mat{A}_I\dot q$, un déplacement virtuel du point
matériel en $I$ s'écrit $\delta\vect{r}_I = \mat{A}_I\,\delta q$ :
$\mat{A}_I$ est la jacobienne $\mat{J}_I$ de ce point. Le travail
virtuel de la force de contact $\vect{N}$ exercée par le sol en $I$ est
donc
\[
    \delta W_r = \vect{N}\cdot\delta\vect{r}_I
    = \vect{N}^T\mat{A}_I\,\delta q
    = (\mat{A}_I^T\vect{N})^T\delta q
\]
En identifiant avec $\vect{Q}^r = \mat{A}_I^T\vect{\lambda}$, on obtient
que $\vect{\lambda}$ est la colonne des coordonnées de la force de contact
$\vect{N}$ dans $\rep{s}$. Ainsi $\lambda_x$ est le frottement
tangentiel. On garde donc la seule contrainte utile, $\mat{A} = (1\ \ -r)$,
de multiplicateur $\lambda_x$ :
\[
    \vect{Q}^r = \mat{A}^T\lambda_x = \begin{pmatrix} \lambda_x\\ -\,r\,\lambda_x
    \end{pmatrix}
\]
où $-\,r\lambda_x$ est le moment du frottement par rapport à $C$. Enfin,
pour tout $\delta q\in\Ker\mat{A}$, c'est-à-dire
$\delta x = r\,\delta\varphi$,
\[
    (\vect{Q}^r)^T\delta q = \lambda_x\,(\delta x - r\,\delta\varphi) = 0
\]
la liaison est donc bien idéale.

\paragraph{Équations du mouvement}

On calcule les termes du lagrangien, il faut se rappeler que $q$ et $\dot q$
sont indépendantes :
\[
    \pd{L}{\dot x} = m\,\dot x,
    \qquad
    \pd{L}{\dot\varphi} = J\,\dot\varphi,
    \qquad
    \pd{L}{x} = -\,k\,x,
    \qquad
    \pd{L}{\varphi} = 0
\]
donc
\[
    \mathcal{L}(L) = \ddt\pd{L}{\dot q} - \pd{L}{q}
    = \begin{pmatrix} m\,\ddot x + k\,x\\ J\,\ddot\varphi \end{pmatrix}
\]
Les équations d'Euler-Lagrange-d'Alembert
$\mathcal{L}(L) = \vect{Q}^{\mathrm{nc}} + \mat{A}^T\lambda_x$ et la
contrainte donnent
\[
    \left\{
    \begin{aligned}
        m\,\ddot x + k\,x &= F + \lambda_x\\
        J\,\ddot\varphi &= -\,r\,\lambda_x\\
        \dot x &= r\,\dot\varphi
    \end{aligned}
    \right.
\]

\paragraph{Résolution}

La contrainte donne $\ddot\varphi = \ddot x/r$, donc
$\lambda_x = -J\ddot x/r^2$. En réinjectant dans la première équation,
\[
    \Bigl(m + \frac{J}{r^2}\Bigr)\ddot x + k\,x = F
\]
C'est un oscillateur harmonique de masse effective
$m_{\mathrm{eff}} = m + J/r^2$, de pulsation et de position d'équilibre :
\[
    \omega_0 = \sqrt{\frac{k}{m_{\mathrm{eff}}}},
    \qquad
    x_{\mathrm{eq}} = \frac{F}{k}
\]
La solution est $x(t) = x_{\mathrm{eq}} + A\cos(\omega_0 t + \phi)$,
d'amplitude $A\geq 0$ et de phase $\phi$ fixées par les conditions
initiales (on ne s'y intéresse pas).

\paragraph{Validité du non-glissement}

Le modèle suppose le roulement sans glissement, mais rien dans les
équations ne borne $\lambda_x$ pour que le modèle valide la modélisation de non glissement. Il faut donc ainsi vérifier la validité du modèle, avec $\lambda_y = m g$ la réaction normale du sol et $\mu$ le
coefficient de frottement :
\[
    |\lambda_x| \leq \mu\,\lambda_y = \mu\,m\,g
\]
Comme $x - x_{\mathrm{eq}} = A\cos(\omega_0 t + \phi)$, on a
$|x - x_{\mathrm{eq}}|\leq A$, avec égalité aux extrémités de
l'oscillation. Le frottement y est maximal :
\[
    |\lambda_x|_{\max} = \frac{J}{r^2}\,\omega_0^2\,A
\]
Pour un disque plein rigide, $J = \tfrac12 m r^2$, donc $J/r^2 = m/2$,
$m_{\mathrm{eff}} = \tfrac32 m$ et $\omega_0^2 = \frac{2k}{3m}$, d'où
\[
    |\lambda_x|_{\max} = \frac{m}{2}\cdot\frac{2k}{3m}\,A = \frac{k A}{3}
\]
La roue roule sans glisser tant que $\frac{kA}{3}\leq\mu\,m\,g$,
c'est-à-dire
\[
    A \leq \frac{3\,\mu\,m\,g}{k}
\]
Au-delà de cette amplitude, la roue glisse ou la contrainte $\dot x = r\dot\varphi$ ne
se tient plus et il faut changer de modèle. Tant qu'elle le tient, le point
de contact a une vitesse nulle et $\lambda_x$ ne travaille pas et la
liaison est bien idéale.
\newpage
